%% kapa: exam=1 round=24 year=2013 period=2 form=A subject=accounting subject_ko=회계학 seq=6 number=6
%% source: Q-Net 감정평가사 기출문제 (cst00309) artlSeq=5073802 file=extracted/1008/5073802/2013년 제24회 감정평가사 1차시험 기출문제/2교시(A형).hwp
%% conversion: kordoc 4.24.51 -> tools/gen1.py (auto); review=proofread-sonnet; flags=none
\begin{question}{6}
(주)한국은 20×1년 7월 1일 사옥신축을 위하여 (주)민국과 건설계약을 체결하였다. 사옥의 건설기간은 20×1년 7월 1일부터 20×2년 12월 31일까지이다. (주)한국은 동 사옥건설과 관련하여 20×1년 7월 1일과 10월 1일에 각각 ￦3,000,000과 ￦6,000,000을 지출하였다. 건설기간동안 (주)한국의 차입금과 관련된 내용은 다음과 같다. \par
{\small\centering
\begin{tabular}{|c|c|c|c|}
\hline
\makecell{차입금} & \makecell{A} & \makecell{B} & \makecell{C} \\
\hline
\makecell{차입금액} & \makecell{￦2,000,000} & \makecell{￦2,000,000} & \makecell{?} \\
\hline
\makecell{차입일} & \makecell{20×0.9.1} & \makecell{20×1.1.1} & \makecell{20×1.7.1} \\
\hline
\makecell{상환일} & \makecell{20×3.8.31} & \makecell{20×2.12.31} & \makecell{20×3.6.30} \\
\hline
\makecell{연이자율} & \makecell{10\%} & \makecell{8\%} & \makecell{10\%} \\
\hline
\makecell{이자지급조건} & \multicolumn{3}{|c|}{\makecell{매년말 후급}} \\
\hline
\end{tabular}\par}

차입금 A와 차입금 B는 일반차입금이며, 차입금 C는 동 사옥건설을 위하여 차입한 특정차입금으로 (주)한국은 차입과 동시에 사옥건설을 위하여 전액 지출하였다. (주)한국이 동 사옥건설과 관련하여 20×1년도에 자본화한 일반차입금의 차입원가가 ￦135,000이라면, (주)한국이 금융기관으로부터 차입한 차입금 C는? (단, 이자는 월할계산한다.) \par
\begin{choices}
  \item ￦750,000
  \item ￦1,000,000
  \item ￦2,000,000
  \item ￦2,500,000
  \item ￦3,000,000
\end{choices}
\end{question}
\answer{5}
